\section*{Chapter \ref{ch:g10:habitable-planet} --- A Habitable Planet}

\begin{solution}{exo:g10:habitable-planet:1}
The temperature at which a perfect emitter would radiate to space exactly the
power the planet absorbs from its star. The Earth's surface is warmer
(\qty{288}{K} against about \qty{255}{K}) because greenhouse gases absorb
part of the infrared it emits and send part of it back down.
\end{solution}

\begin{solution}{exo:g10:habitable-planet:2}
Venus: measured \qty{737}{K} against an \omterm{def:g10:habitable-planet:equilibrium-temperature}{equilibrium temperature} of about
\qty{227}{K}, a warming of about \qty{510}{K}.
\end{solution}

\begin{solution}{exo:g10:habitable-planet:3}
$S = 1361/1.52^2 \approx \qty{589}{W/m^2}$ (the fact sheet gives
\qty{586}{W/m^2}, with the exact mean distance).
\end{solution}

\begin{solution}{exo:g10:habitable-planet:4}
\qty{1}{bar}, \qty{330}{K}: liquid (below the boiling point, about
\qty{373}{K}). \qty{0.001}{bar}, \qty{250}{K}: ice (the pressure lies above
the ice--vapour curve, which is at about \qty{0.0008}{bar} at \qty{250}{K}).
\qty{100}{bar}, \qty{700}{K}: vapour --- the temperature is above the critical
temperature, so no liquid can exist.
\end{solution}

\begin{solution}{exo:g10:habitable-planet:5}
A planet orbiting a star other than the Sun. Transits: the planet passes in
front of its star and dims it slightly at every orbit. Wobble: the planet's
pull moves the star back and forth along our line of sight, which shifts the
colour of its light periodically.
\end{solution}

\begin{solution}{exo:g10:habitable-planet:6}
$T = (1361 \times 0.706/(4 \times 5.67\times10^{-8}))^{1/4} \approx
\qty{255}{K} \approx \qty{-18}{\degreeCelsius}$. With $A = 0$:
$\approx \qty{278}{K} \approx \qty{5}{\degreeCelsius}$.
\end{solution}

\begin{solution}{exo:g10:habitable-planet:7}
$\sqrt 4 = 2$: from $2 \times 0.99 = \qty{1.98}{au}$ to $2 \times 1.70 =
\qty{3.40}{au}$. At \qty{1}{au} the planet would receive four times the
Earth's irradiance: far inside the inner edge, too hot.
\end{solution}

\begin{solution}{exo:g10:habitable-planet:8}
$S = 1361/0.723^2 \approx \qty{2604}{W/m^2}$;
$T_{\mathrm{eq}} = (2604 \times 0.23/(4\sigma))^{1/4} \approx \qty{227}{K}$.
The surface is about \qty{510}{K} hotter: its thick carbon-dioxide atmosphere
(\qty{92}{bar}) produces an enormous greenhouse effect.
\end{solution}

\begin{solution}{exo:g10:habitable-planet:9}
At \qty{6.36}{mbar} water is liquid only between about \qty{0}{\degreeCelsius}
and \qty{0.5}{\degreeCelsius}. At \qty{10}{\degreeCelsius} the pressure is
below the boiling curve: water boils away into vapour; at night it freezes.
\end{solution}

\begin{solution}{exo:g10:habitable-planet:10}
$11.2/2.4 \approx 4.7$. A gas molecule must move almost five times faster to
escape the Earth than the Moon; over billions of years the Moon's weak gravity
let its gases escape, while the Earth's kept most of them.
\end{solution}

\begin{solution}{exo:g10:habitable-planet:11}
$T = (1361 \times 0.50/(4\sigma))^{1/4} \approx \qty{234}{K}$, about
\qty{-39}{\degreeCelsius}. The ice would reflect so much sunlight that the
planet would stay cold enough to keep it: a reinforcing loop.
\end{solution}

\begin{solution}{exo:g10:habitable-planet:12}
$T_{\mathrm{eq}} = (S_0(1-A)/(4\sigma d^2))^{1/4} = (S_0/(4\sigma))^{1/4}
(1-A)^{1/4} (d^2)^{-1/4}$; $(1361/(4\sigma))^{1/4} \approx \qty{278.3}{K}$ and
$(d^2)^{1/4} = \sqrt d$. Four times farther: $T$ divided by $\sqrt 4 = 2$.
\end{solution}

\begin{solution}{exo:g10:habitable-planet:13}
$\sqrt d = 278.3 \times 0.706^{1/4}/273 \approx 0.934$, so $d \approx
\qty{0.873}{au}$, inside the inner edge (\qty{0.99}{au}). The \omterm{def:g10:habitable-planet:equilibrium-temperature}{equilibrium
temperature} ignores the greenhouse effect: a real planet with an atmosphere
is warmer than $T_{\mathrm{eq}}$, so it overheats and loses its water already
farther out, at about \qty{0.99}{au}. The edges come from models of the whole
atmosphere, not from $T_{\mathrm{eq}}$ alone.
\end{solution}

\begin{solution}{exo:g10:habitable-planet:14}
Between the melting point, about \qty{0}{\degreeCelsius}, and the boiling
point, about \qty{7}{\degreeCelsius}. At \qty{15}{\degreeCelsius} the lake
would boil.
\end{solution}

\begin{solution}{exo:g10:habitable-planet:15}
$\sqrt{0.5} \approx 0.707$: from $0.707 \times 0.99 \approx \qty{0.70}{au}$
to $0.707 \times 1.70 \approx \qty{1.20}{au}$. The planet is right at the
outer edge. The edges are model results and the scaling is approximate, so it
is impossible to say whether it is inside or outside.
\end{solution}

\begin{solution}{pb:g10:habitable-planet:1}
\textbf{1.} $0.00151/0.0485^2 \approx 0.642$.

\textbf{2.} $0.642 \times 1361 \approx \qty{874}{W/m^2}$.

\textbf{3.} Between the two: about three times less than Venus, one and a
half times more than Mars.

\textbf{4.} $365.25/11.18 \approx 32.7$.

\textbf{5.} $(874/(4 \times 5.67\times10^{-8}))^{1/4} \approx \qty{249}{K}$.

\textbf{6.} $(874 \times 0.23/(4\sigma))^{1/4} \approx \qty{173}{K}$.

\textbf{7.} Neither: \qty{249}{K} is still below \qty{273}{K}. A bright cover
of clouds or ice lowers the \omterm{def:g10:habitable-planet:equilibrium-temperature}{equilibrium temperature} by more than
\qty{75}{K}: what covers a planet matters as much as its distance.

\textbf{8.} $\sqrt{0.00151} \approx 0.0389$.

\textbf{9.} $0.99 \times 0.0389 \approx \qty{0.0385}{au}$.

\textbf{10.} $1.70 \times 0.0389 \approx \qty{0.066}{au}$.

\textbf{11.} Yes: $0.0385 < 0.0485 < 0.066$.

\textbf{12.} $d_\odot = 0.0485/0.0389 \approx \qty{1.25}{au}$: in the inner
half of the Sun's \omterm{def:g10:habitable-planet:habitable-zone}{habitable zone}, between the Earth and Mars.

\textbf{13.} $1.055^{1/3} \approx 1.018$.

\textbf{14.} $11.2 \times \sqrt{1.055/1.018} \approx \qty{11.4}{km/s}$, about the
Earth's: its gravity could hold an atmosphere (whether it has one depends on
its history).

\textbf{15.} The wobble of the star reveals only the part of its motion along
our line of sight; since the tilt of the orbit is unknown, the mass found is
the true mass only if the orbit is seen edge-on, and a minimum otherwise.

\textbf{16.} The pressure of the triple point, \qty{611.657}{Pa}, about
\qty{0.006}{bar}.

\textbf{17.} $(874 \times 0.706/(4\sigma))^{1/4} \approx \qty{228}{K}$.

\textbf{18.} $288 - 228 = \qty{60}{K}$: Proxima~b would need a greenhouse
warming of \textbf{about \qty{60}{K}}, almost twice the Earth's \qty{33}{K}.
\end{solution}
